\begin{helper}
For every fixed real number \(c>0\),
\[
\frac{\log\bigl(c\,\noderef{RamseyScale}(k)\bigr)}{\log k}\to 2
\]
as \(k\to\infty\).
\end{helper}
\begin{proof}
Fix \(c>0\).  For all sufficiently large \(k\), we have \(k>1\), hence
\(\log k>0\), and therefore \(\noderef{RamseyScale}(k)=k^2/\log k>0\).  Since
\(c>0\), all logarithms below are then taken of positive quantities.

Using the definition of \(\noderef{RamseyScale}\), the logarithm laws for
positive arguments give, eventually in \(k\),
\[
\begin{aligned}
\log\bigl(c\,\noderef{RamseyScale}(k)\bigr)
&=\log c+\log\left(\frac{k^2}{\log k}\right)\\
&=\log c+\log(k^2)-\log(\log k)\\
&=\log c+2\log k-\log\log k.
\end{aligned}
\]
Dividing by \(\log k>0\), we obtain the eventual identity
\[
\frac{\log\bigl(c\,\noderef{RamseyScale}(k)\bigr)}{\log k}
=2+\frac{\log c}{\log k}-\frac{\log\log k}{\log k}.
\]
The first error term tends to \(0\), because \(\log c\) is a fixed real number
and \(\log k\to\infty\).

For the second error term, put \(t=\log k\).  Since \(k\to\infty\), we have
\(t\to\infty\).  Once \(t\ge1\), the elementary estimate
\(\log t\le\sqrt t\) holds.  Hence eventually
\[
0\le\frac{\log\log k}{\log k}
=\frac{\log t}{t}
\le\frac{\sqrt t}{t}
=\frac1{\sqrt t}.
\]
The right hand side tends to \(0\) because \(t\to\infty\), so by squeezing
\[
\frac{\log\log k}{\log k}\to0.
\]
Both error terms in the eventual identity tend to \(0\), and therefore
\[
\frac{\log\bigl(c\,\noderef{RamseyScale}(k)\bigr)}{\log k}\to2 .
\]
\end{proof}
